\SourceLeaf{31}\[
z\equiv\left\{\begin{array}{rl}
17 & (\mathrm{mod.}\ 9)\\
-4 & (\mathrm{mod.}\ 5)\\
-4 & (\mathrm{mod.}\ 7)\\
33 & (\mathrm{mod.}\ 16)
\end{array}\right.
\quad\text{unde colligitur}\quad z\equiv 3041\ (\mathrm{mod.}\ 5040).
\]

Ceterum palam est, plerumque commodius fore, si de conditionibus remanentibus eae quae
ex una eademque conditione evolutae erant seorsim recolligantur, quum hoc nullo negotio
fieri possit; \emph{e.g.} quando ex conditionibus $z\equiv a\ (\mathrm{mod.}\ A')$,
$z\equiv a\ (\mathrm{mod.}\ A'')$ etc. aliquae abierunt: quae ex reliquis restituitur,
haec erit, $z\equiv a$ secundum modulum qui est productum omnium modulorum ex
$A',A'',A'''$ etc. remanentium. Ita in nostro exemplo ex conditionibus
$z\equiv -4\ (\mathrm{mod.}\ 5)$, $z\equiv -4\ (\mathrm{mod.}\ 7)$ ea ex qua ortae erant
$z\equiv -4\ (\mathrm{mod.}\ 35)$ sponte restituitur. Porro hinc sequitur haud prorsus
perinde esse, quaenam ex conditionibus superfluis reiiciantur, quantum ad calculi
brevitatem: sed haec aliaque artificia practica, quae ex usu multo facilius quam ex
praeceptis ediscuntur hic tradere non est instituti nostri.

\Article{36}

Quando omnes moduli $A,B,C,D$ etc. inter se sunt primi, sequenti methodo saepius
praestat uti. Determinetur numerus $\alpha$ secundum $A$ unitati, secundum reliquorum
modulorum productum vero cifrae congruus, sive sit $\alpha$ valor quicunque (plerumque
praestat \emph{minimum} accipere) expressionis $\frac{1}{BCD\ \mathrm{etc.}}\ (\mathrm{mod.}\ A)$
per $BCD$ etc. multiplicatus (vid.\ art.\ 32); similiter sit
$\beta\equiv 1\ (\mathrm{mod.}\ B)$ et $\equiv 0\ (\mathrm{mod.}\ ACD\ \mathrm{etc.})$,
$\gamma\equiv 1\ (\mathrm{mod.}\ C)$ et $\equiv 0\ (\mathrm{mod.}\ ABD\ \mathrm{etc.})$, etc.
Tunc si numerus $z$ desideratur, qui secundum modulos $A,B,C,D$ etc. numeris
$a,b,c,d$ etc. respective sit congruus, poni poterit

\[
z\equiv\alpha a+\beta b+\gamma c+\delta d\ \mathrm{etc.}\ (\mathrm{mod.}\ ABCD\ \mathrm{etc.}).
\]

Manifesto enim, $\alpha a\equiv a\ (\mathrm{mod.}\ A.)$; reliqua autem membra
$\beta b,\gamma c$ etc. omnia $\equiv 0\ (\mathrm{mod.}\ A)$: quare
$z\equiv a\ (\mathrm{mod.}\ A)$. Similiter de reliquis modulis demonstratio adornatur.
Haec solutio priori praeferenda, quando plura huiusmodi problemata sunt solvenda, pro
quibus moduli $A,B,C$ etc. valores suos retinent; tunc enim numeri $\alpha,\beta,\gamma$
etc., valores constantes nanciscuntur. Hoc usu venit in problemate chronologico ubi
quaeritur, quotus in periodo Juliana sit annus, cuius indictio, numerus aureus, et
cyclus solaris dantur. Hic $A=15$, $B=19$, $C=28$; quare,
\SourceLeaf{32}quum valor expressionis $\frac{1}{19.28}\ (\mathrm{mod.}\ 15)$, sive $\frac{1}{532}\ (\mathrm{mod.}\ 15)$, sit $13$, erit $\alpha=6916$. Similiter pro $\beta$ invenitur $4200$, et pro $\gamma$ $4845$, quare numerus quaesitus erit residuum minimum numeri $6916a+4200b+4845c$, denotantibus $a$ indictionem. $b$ numerum aureum, $c$ cyclum solarem.

\begin{center}\emph{Congruentiae lineares quae plures incognitas implicant.}\end{center}

\Article{37}

Haec de congruentiis primi gradus unicam incognitam continentibus sufficiant. Superest ut de congruentiis agamus, in quibus plures incognitae sunt permixtae. At quoniam hoc caput, si omni rigore singula exponere velimus, sine prolixitate absolvi non potest, propositumque hoc loco nobis non est, omnia exhaurire, sed ea tantum tradere, quae attentione digniora videantur: hic ad paucas observationes investigationem restringimus, uberiorem huius rei expositionem ad aliam occasionem nobis reservantes.

1) Simili modo, ut in aequationibus, perspicitur, etiam hic totidem congruentias haberi debere, quot sint incognitae determinandae.

2) Propositae sint igitur congruentiae

\[
\begin{array}{l@{\qquad}c@{\qquad}r}
ax+by+cz\ldots\equiv f\ (\mathrm{mod.}\ m)&\ldotp\quad\ldotp\quad\ldotp\quad\ldotp\quad\ldotp&(A)\\
a'x+b'y+c'z\ldots\equiv f'&\ldotp\quad\ldotp\quad\ldotp\quad\ldotp\quad\ldotp&(A')\\
a''x+b''y+c''z\ldots\equiv f''&\ldotp\quad\ldotp\quad\ldotp\quad\ldotp\quad\ldotp\quad\ldotp\quad\ldotp&(A'')\\
\mathrm{etc.}&&
\end{array}
\]

 totidem numero, quot sunt incognitae $x,y,z$ etc.

Iam determinentur numeri $\xi,\xi',\xi''$ etc. ita ut sit

\[\begin{array}{l}
b\xi+b'\xi'+b''\xi''+\mathrm{etc.}=0\\
c\xi+c'\xi'+c''\xi''+\mathrm{etc.}=0\\
\qquad\mathrm{etc.}
\end{array}\]

 et quidem ita ut omnes sint integri nullumque factorem communem habeant, quod fieri posse ex theoria aequationum linearium constat. Simili modo determinentur $v,v',v''$ etc., $\zeta,\zeta',\zeta''$ etc. etc. ita ut sit

\[
\begin{array}{l}av+a'v'+a''v''+\mathrm{etc.}=0\\cv+c'v'+c''v''+\mathrm{etc.}=0\\\qquad\mathrm{etc.}\end{array}
\]


\SourceLeaf{33}\[\begin{array}{l}
a\zeta+a'\zeta'+a''\zeta''+\mathrm{etc.}=0\\
b\zeta+b'\zeta'+b''\zeta''+\mathrm{etc.}=0\\
\qquad\mathrm{etc.}\quad\mathrm{etc.}
\end{array}\]

3) Manifestum est si congruentiae $A,A',A''$ etc. per $\xi,\xi',\xi''$ etc., tum per $v,v',v'',$ etc. etc. multiplicentur. tuncque addantur, has congruentias proventuras esse:

\[
\begin{array}{l}(a\xi+a'\xi'+a''\xi''+\mathrm{etc.})x\equiv f\xi+f'\xi'+f''\xi''+\mathrm{etc.}\\(bv+b'v'+b''v''+\mathrm{etc.})y\equiv fv+f'v'+f''v''+\mathrm{etc.}\\(c\zeta+c'\zeta'+c''\zeta''+\mathrm{etc.})z\equiv f\zeta+f'\zeta'+f''\zeta''+\mathrm{etc.}\\\qquad\mathrm{etc.}\end{array}
\]

 quas brevitatis gratia ita exhibemus:

\[
\Sigma(a\xi)x\equiv\Sigma(f\xi),\quad\Sigma(bv)y\equiv\Sigma(fv),\quad\Sigma(c\zeta)z\equiv\Sigma(f\zeta)\ \mathrm{etc.}
\]

4) Iam plures casus sunt distinguendi.

\emph{Primo} quando omnes incognitarum coefficientes $\Sigma(a\xi),\Sigma(av)$ etc. ad congruentiarum modulum $m$ sunt primi, hae congruentiae secundum praecepta ante tradita solvi possunt, problematisque solutio completa per congruentias formae $x\equiv p\ (\mathrm{mod.}\ m)$, $y\equiv q\ (\mathrm{mod.}\ m)$ etc. exhibebitur\begingroup\renewcommand{\thefootnote}{*)}\footnote{Observare convenit hancce conclusionem demonstratione egere, quam autem hic supprimimus. Proprie enim nihil aliud ex analysi nostra sequitur, quam quod congruentiae propositae per alios incognitarum $x,y$ etc. valores solvi nequeant: hos vero satisfacere non sequitur. Fieri enim posset ut nulla omnino solutio daretur. Similis paralogismus etiam in aequationum linearium explicatione plerumque committitur.}\endgroup. \emph{E.g.} Si proponuntur congruentiae

\[x+3y+z\equiv1,\quad4x+y+5z\equiv7,\quad2x+2y+z\equiv3\ (\mathrm{mod.}\ 8)\]

 invenietur $\xi=9$, $\xi'=1$, $\xi''=-14$, unde fit $-15x\equiv-26$, quare $x\equiv6\ (\mathrm{mod.}\ 8)$; eodem modo invenitur $15y\equiv-4$, $15z\equiv1$, et hinc $y\equiv4$, $z\equiv7\ (\mathrm{mod.}\ 8)$.

5) \emph{Secundo} quando non omnes coefficientes $\Sigma(a\xi),\Sigma(bv)$ etc. ad modulum sunt primi, sint $\alpha,\beta,\gamma$ etc. divisores communes maximi ipsius $m$ cum $\Sigma(a\xi),\Sigma(bv),\Sigma(c\zeta)$ etc. resp., patetque problema impossibile esse, nisi illi numeros $\Sigma(f\xi),\Sigma(fv),\Sigma(f\zeta)$ etc. resp. metiantur. Quando vero hae conditiones locum habent, congruentiae in (3) complete resolventur per tales $x\equiv p\ (\mathrm{mod.}\ \frac{m}{\alpha})$, $y\equiv q\ (\mathrm{mod.}\ \frac{m}{\beta})$, $z\equiv r\ (\mathrm{mod.}\ \frac{m}{\gamma})$ etc., aut si mavis dabuntur $\alpha$ valores diversi ipsius $x$ (i. e. secundum $m$ incongrui puta $p,p+\frac{m}{\alpha}\ldots p+\frac{(\alpha-1)m}{\alpha}$),
\SourceLeaf{34}$\beta$ valores diversi ipsius $y$ etc., illis congruentiis satisfacientes: manifestoque omnes solutiones congruentiarum propositarum (si quae omnino dantur) inter illas reperientur. Attamen hanc conclusionem convertere non licet; nam plerumque non omnes combinationes omnium $\alpha$ valorum ipsius $x$ cum omnibus ipsius $y$ cum omnibus ipsius $z$ etc. problemati satisfaciunt, sed quaedam tantum, quarum nexum per unam pluresve congruentias conditionales exhibere licet. At quum completa huius problematis resolutio ad sequentia non sit necessaria, hoc argumentum fusius hoc loco non exsequimur, exemploque ideam qualemcunque de eo dedisse sat habemus.

Propositae sint congruentiae

\[3x+5y+z\equiv4,\quad2x+3y+2z\equiv7,\quad5x+y+3z\equiv6\ (\mathrm{mod.}\ 12)\]

 Hic fiunt $\xi,\xi',\xi''$; $v,v',v''$; $\zeta,\zeta',\zeta''$; resp. $=1,-2,1$; $1,1,-1$; $-13,22,-1$, unde $4x\equiv-4$, $7y\equiv5$, $28z\equiv96$. Hinc prodeunt quatuor valores ipsius $x$ puta $\equiv2,5,8.\ 11$; unus valor ipsius $y$ puta $\equiv11$; quatuor valores ipsius $z$ puta $\equiv0,3,6,9\ (\mathrm{mod.}\ 12)$. Iam ut sciamus, quasnam combinationes valorum ipsius $x$ cum valoribus ipsius $z$ adhibere liceat, substituimus in congruentiis propp. pro $x,y,z$ resp. $2+3t$, $11$, $3u$, unde transeunt in has

\[57+9t+3u\equiv0,\quad30+6t+6u\equiv0,\quad15+15t+9u\equiv0\ (\mathrm{mod.}\ 12)\]

 quibus facile intelligitur aequivalere has

\[19+3t+u\equiv0,\quad10+2t+2u\equiv0,\quad5+5t+3u\equiv0\ (\mathrm{mod.}\ 4)\]

 Prima manifesto requirit ut sit $u\equiv t+1\ (\mathrm{mod.}\ 4)$, quo valore in reliquis substituto etiam his satisfieri invenitur. Hinc colligitur, valores ipsius $x$ hos $2,5,8,11$ (qui prodeunt statuendo $t\equiv0.\ 1,2,3$) necessario combinandos esse cum valoribus ipsius $z$ his $z\equiv3,6,9,0$ resp., ita ut omnino quatuor solutiones habeantur

\[\begin{array}{r@{\;}r@{,\quad}r@{,\quad}r@{,\quad}r@{\ }l}
x\equiv&2&5&8&11&(\mathrm{mod.}\ 12)\\
y\equiv&11&11&11&11&\\
z\equiv&3&6&9&0&
\end{array}\]

\begin{center}$\ast\quad\underset{\ast}{\phantom{\ast}}\quad\ast$\end{center}
